\section{The customer bought capability, not coaching}

The public encounter with a frontier model increasingly begins with a seemingly simple bargain: describe a task and receive useful work. The user may be able to see that the model understands the request, drafts a competent plan, and names the next tool action. Yet the consumer product stops at a caveat, a draft, a request for another prompt, or a partially completed sequence. The user then becomes the missing workflow engine: reframing the request, collecting context, moving information between applications, restarting after an interruption, and deciding whether the external state actually changed. A product that can explain the work but cannot finish it has sold less than its demonstration suggested.

This paper calls the missing layer an \emph{aftermarket harness}: a specialist system added around a consumer-accessible frontier product or model to make a narrow class of legitimate tasks complete reliably. The harness may transform an open request into a testable task, supply permitted context and tools, manage long-lived state, route an exception to a person, and verify the result. Those functions are not synonymous with a longer prompt. They are software and research choices about what the model sees, what it can do, when an attempt should stop, and what counts as an accepted outcome. Work on compound AI systems and agent-computer interfaces supports treating this surrounding system as a first-class object of design \citep{compound,sweagent}. It does not establish that any particular aftermarket business is profitable.

Our commercial proposition is deliberately provocative: for some teams outside frontier laboratories, the highest-leverage research agenda may be to unlock work that capable consumer AI products currently decline or leave unfinished without repeated human nudging. ``Highest-leverage'' is a comparative claim. It must be tested against better model access, a frontier provider's own product improvements, conventional automation, and other uses of an applied research budget. We do not infer it from a benchmark score, an isolated demo, or the existence of a user complaint. Its appeal is that it asks a concrete question: can a specialist vendor produce more verified, authorized work from the same underlying intelligence at lower total cost?

The product, rather than the bare model, is the relevant starting point. A consumer AI product includes a model, instructions, routing, tools, context window, budgets, interfaces, policies, account connections, and handoff rules. Observing a decline or incomplete result tells us what that product did in that episode; it does not reveal the model's maximum latent ability. Conversely, a model that can be coaxed into generating a plausible plan has not demonstrated the ability to execute the plan. Figure~\ref{fig:wedge} represents the candidate gap, but the curves are only a research hypothesis. The region of commercial interest exists only for tasks where the harness improves verified completion through a permitted route.

\begin{figure}[tbp]
\centering\begin{tikzpicture}
\begin{axis}[
  width=\linewidth,height=\ifdefined\SlidePlotHeight\SlidePlotHeight\else 6.2cm\fi,
  xmin=0,xmax=10,ymin=0,ymax=1,
  axis lines=left,xtick=\empty,ytick=\empty,
  xlabel={\ifdefined\SlidePlotHeight Task demands\else Composite task demands $\longrightarrow$\fi},
  ylabel={\ifdefined\SlidePlotHeight Completion\else Authorized completion propensity $\longrightarrow$\fi},
  xlabel style={font=\small},ylabel style={font=\small},
  axis line style={ink!65},clip=false]
\addplot[name path=base,draw=native,line width=1.6pt,smooth] coordinates {(0,.87)(2,.77)(4,.55)(6,.29)(8,.13)(10,.05)};
\addplot[name path=nudge,draw=ochre,dashed,line width=1.6pt,smooth] coordinates {(0,.90)(2,.84)(4,.71)(6,.52)(8,.28)(10,.12)};
\addplot[name path=harness,draw=teal,line width=2pt,smooth] coordinates {(0,.91)(2,.87)(4,.79)(6,.68)(8,.48)(10,.25)};
\addplot[teal!14] fill between[of=harness and base];
\node[font=\footnotesize,text=native,anchor=west] at (axis cs:6.55,.22) {Unassisted};
\node[font=\footnotesize,text=ochre,anchor=west] at (axis cs:6.1,.49) {Nudged};
\node[font=\footnotesize,text=teal,anchor=west] at (axis cs:6.2,.73) {Specialist harness};
\node[font=\footnotesize,text=ink,align=center] at (axis cs:3.7,.39) {Hypothesized\\completion wedge};
\end{axis}
\end{tikzpicture}

\caption{A schematic completion wedge between a stock consumer product, manual nudging, and a specialist harness. The curves and shaded area are conceptual, not empirical probabilities or market shares. Some stops reflect real incapacity or legitimate constraints and cannot be recovered by a harness.}
\label{fig:wedge}
\end{figure}

The thesis differs from the familiar claim that models will automate many tasks. It identifies a \emph{research location}: between accessible model competence and the customer's accepted result. This location is attractive to a non-frontier team because it can be studied with a fixed model. Hold the model version and task constant; change task specification, tool interface, memory, control loop, or verification; then compare accepted outcomes and total cost. A gain under that comparison is evidence about the harness. It is not proof that the model was secretly capable of every task it refused, nor that the same intervention will generalize to another service.

There is also a market timing argument. Frontier providers are themselves improving tools and agent harnesses. Anthropic's published engineering accounts describe substantial differences between a generic loop and a designed long-running workflow \citep{anthropiclong,anthropic2026}. Those accounts are evidence that harness design matters within particular vendor experiments, and they are a warning that an aftermarket advantage may be temporary. A specialist must find tasks for which domain context, service integration, acceptance criteria, or customer trust evolves faster than a broad provider product can absorb. The right experiment therefore tests both current uplift and how much of it survives a better native product.

The unit sold is neither a token nor an impressive conversation. It is a task whose requested external state has been independently checked, with the owner’s required approvals accounted for. Once the unit is stated this way, commercially meaningful questions become possible: How often does the stock product finish? How many minutes of expert nudging does the user contribute? Which harness intervention changes the outcome? How much does verification and maintenance cost? We take these questions, rather than a universal claim of model reluctance, as the paper's starting point.
